% alt: a chain of twenty-two sites with a region R of nine sites holding a core R_0 of three sites in its middle, a buffer Q of two sites on either side of it (two at the left end of the chain, two between R and C), a region C of nine sites holding a core C_0 of six sites at the right end, three brackets of width w under the three sites separating each core from the edge of its region, and two shaded bands marking the support S of the resetting unitary, which covers the buffer and the cut edges of R, reaches two sites into each margin on the R side and one site into the margin on the C side, and leaves both cores and the innermost margin sites clear.
% Section 8.2 (Information reset), Lemma lem:reset and the support step of its proof: disjoint regions R and C with cores R_0 and C_0 at distance more than w = C_p(1+b+log L) from the complements of their regions, and Q = Lambda \ (R cup C) the rest of the chain (drawn for l = R = 1 and w = 3); the dots are the sites and the upward legs the physical indices C^q, every leg is open (nothing is contracted); the red enclosures are the cores R_0 (sites 6 to 8) and C_0 (sites 17 to 22), the blue enclosure is R (sites 3 to 11), the purple ones are the two pieces of Q (sites 1 to 2 and 12 to 13), the dashed enclosure is C (sites 14 to 22); the three blue brackets under the chain mark the margins of width w = 3 between R_0 and Lambda \ R (sites 3 to 5 and 9 to 11) and between C_0 and Lambda \ C (sites 14 to 16); the two shaded enclosures mark S = N_l(Q cup Z_R), the support of the unitary U: Z_R = {2,3,11,12} is the union of the supports of the terms of H meeting both R and Lambda \ R, so Q cup Z_R = {1,2,3} cup {11,12,13} and its 1-neighbourhood S = {1,...,4} cup {10,...,14} reaches l + R = 2 sites into R at each cut edge (sites 3, 4 and 10, 11, two of the three margin sites) and l = 1 site into C (site 14, one of the three margin sites), so it stays clear of both cores, with margin sites 5, 9 and 15, 16 untouched, because w = 3 exceeds l + R = 2.
\begin{tenkz}[lattice={1x22}, frame=plane, physical=up, bonds=none]
  \tn{} & \tn{} & \tn{} & \tn{} & \tn{} & \tn{} & \tn{} & \tn{} & \tn{} & \tn{} & \tn{} & \tn{} & \tn{} & \tn{} & \tn{} & \tn{} & \tn{} & \tn{} & \tn{} & \tn{} & \tn{} & \tn{}
  \tnmark[form=enclosure, species=neutral, tint, label pos=n]{(1,1) .. (1,4)}{$S$}
  \tnmark[form=enclosure, species=neutral, tint, label pos=n]{(1,10) .. (1,14)}{$S=N_\ell(Q\cup Z_R)$}
  \tnmark[form=enclosure, species=selected, label pos=s]{(1,6) .. (1,8)}{$R_0$}
  \tnmark[form=enclosure, species=secondary, label pos=n]{(1,3) .. (1,11)}{$R$}
  \tnmark[form=enclosure, species=collar, label pos=s]{(1,1) .. (1,2)}{$Q$}
  \tnmark[form=enclosure, species=collar, label pos=s]{(1,12) .. (1,13)}{$Q$}
  \tnmark[form=enclosure, species=complement, label pos=n]{(1,14) .. (1,22)}{$C$}
  \tnmark[form=enclosure, species=selected, label pos=s]{(1,17) .. (1,22)}{$C_0$}
  \tnmark[form=bracket, species=secondary, label pos=s]{(1,3) .. (1,5)}{$w$}
  \tnmark[form=bracket, species=secondary, label pos=s]{(1,9) .. (1,11)}{$w$}
  \tnmark[form=bracket, species=secondary, label pos=s]{(1,14) .. (1,16)}{$w$}
\end{tenkz}
